\section{Evaluation}
\label{sec:Evaluation}

This section presents a comprehensive empirical evaluation of our dynamic Multi-Armed Bandit (MAB) hyper-heuristic framework evaluated against real-world HPC workload traces from HPC2N and NASA Ames iPSC/860. We examine policy evolution dynamics, cumulative makespan trajectories, multi-objective trade-off profiles, and cross-algorithm selection heatmaps.

\subsection{Experimental Results and Comparative Analysis}

\subsubsection{Heuristic Selection Ratio Evolution}
Figures~\ref{fig:selection_evolution_hpc2n} and~\ref{fig:selection_evolution_nasa} track the temporal evolution of low-level heuristic (LLH) selection ratios across 100 simulation rounds using a 10-round moving window.

\begin{figure}[htbp]
    \centering
    \begin{subfigure}[b]{0.48\textwidth}
        \centering
        \includegraphics[width=\textwidth]{hpc2n_ucb_arm_selection_evolution.png}
        \caption{HPC2N -- UCB}
        \label{fig:selection_hpc2n_ucb}
    \end{subfigure}
    \hfill
    \begin{subfigure}[b]{0.48\textwidth}
        \centering
        \includegraphics[width=\textwidth]{hpc2n_discounted_ucb_arm_selection_evolution.png}
        \caption{HPC2N -- Discounted UCB}
        \label{fig:selection_hpc2n_discounted}
    \end{subfigure}
    \vskip\baselineskip
    \begin{subfigure}[b]{0.48\textwidth}
        \centering
        \includegraphics[width=\textwidth]{hpc2n_thompson_arm_selection_evolution.png}
        \caption{HPC2N -- Thompson Sampling}
        \label{fig:selection_hpc2n_thompson}
    \end{subfigure}
    \hfill
    \begin{subfigure}[b]{0.48\textwidth}
        \centering
        \includegraphics[width=\textwidth]{hpc2n_epsilon_arm_selection_evolution.png}
        \caption{HPC2N -- $\epsilon$-Greedy}
        \label{fig:selection_hpc2n_epsilon}
    \end{subfigure}
    \caption{Low-level heuristic selection ratio evolution over 100 rounds on the HPC2N trace.}
    \label{fig:selection_evolution_hpc2n}
\end{figure}

\begin{figure}[htbp]
    \centering
    \begin{subfigure}[b]{0.48\textwidth}
        \centering
        \includegraphics[width=\textwidth]{nasa_ucb_arm_selection_evolution.png}
        \caption{NASA -- UCB}
        \label{fig:selection_nasa_ucb}
    \end{subfigure}
    \hfill
    \begin{subfigure}[b]{0.48\textwidth}
        \centering
        \includegraphics[width=\textwidth]{nasa_discounted_ucb_arm_selection_evolution.png}
        \caption{NASA -- Discounted UCB}
        \label{fig:selection_nasa_discounted}
    \end{subfigure}
    \vskip\baselineskip
    \begin{subfigure}[b]{0.48\textwidth}
        \centering
        \includegraphics[width=\textwidth]{nasa_thompson_arm_selection_evolution.png}
        \caption{NASA -- Thompson Sampling}
        \label{fig:selection_nasa_epsilon}
    \end{subfigure}
    \hfill
    \begin{subfigure}[b]{0.48\textwidth}
        \centering
        \includegraphics[width=\textwidth]{nasa_epsilon_arm_selection_evolution.png}
        \caption{NASA -- $\epsilon$-Greedy}
        \label{fig:selection_hpc2n_epsilon}
    \end{subfigure}
    \caption{Low-level heuristic selection ratio evolution over 100 rounds on the NASA Ames trace.}
    \label{fig:selection_evolution_nasa}
\end{figure}

\begin{itemize}
    \item \textbf{Initial Exploration Phase ($k \le 10$):} Unbiased initial sampling yields a broad distribution over the selected algorithms
    \item \textbf{Exploitation Regime on HPC2N ($k > 15$):} On the HPC2N trace, UCB (Figure~\ref{fig:selection_hpc2n_ucb}) and Discounted UCB (Figure~\ref{fig:selection_hpc2n_discounted}) heavily exploit \texttt{hghhc}, which occupies up to $80\%$ of selection windows between rounds 40 and 90, while selectively triggering \texttt{heft} and \texttt{peft} during workload shifts around round 65 and 95. Thompson Sampling (Figure~\ref{fig:selection_hpc2n_thompson}) maintains a more continuous sampling spread between \texttt{hghhc}, \texttt{peft}, and \texttt{sufferage}.
    \item \textbf{Phase Lock on NASA Ames:} On the NASA Ames trace, UCB (Figure~\ref{fig:selection_nasa_ucb}) locks onto \texttt{round\_robin} and \texttt{max\_min} between rounds 30 and 60 before transitioning to \texttt{peft}. Discounted UCB (Figure~\ref{fig:selection_nasa_discounted}) dynamically shuffles between \texttt{sufferage}, \texttt{simulated\_annealing}, \texttt{hghhc}, and \texttt{peft}, which prevents a premature policy saturation in highly uniform arrival regimes.
\end{itemize}

\subsubsection{Cumulative Makespan Trajectories}
Figures~\ref{fig:makespan_hpc2n_all} and~\ref{fig:makespan_nasa_all} showcase the cumulative makespan curves across all four bandit implementations against standalone baseline heuristics.

\begin{figure}[htbp]
    \centering
    \begin{subfigure}[b]{0.48\textwidth}
        \centering
        \includegraphics[width=\textwidth]{hpc2n_ucb_cumulative_makespan.png}
        \caption{HPC2N -- UCB}
        \label{fig:makespan_hpc2n_ucb}
    \end{subfigure}
    \hfill
    \begin{subfigure}[b]{0.48\textwidth}
        \centering
        \includegraphics[width=\textwidth]{hpc2n_discounted_ucb_cumulative_makespan.png}
        \caption{HPC2N -- Discounted UCB}
        \label{fig:makespan_hpc2n_discounted}
    \end{subfigure}
    \vskip\baselineskip
    \begin{subfigure}[b]{0.48\textwidth}
        \centering
        \includegraphics[width=\textwidth]{hpc2n_thompson_cumulative_makespan.png}
        \caption{HPC2N -- Thompson Sampling}
        \label{fig:makespan_hpc2n_thompson}
    \end{subfigure}
    \hfill
    \begin{subfigure}[b]{0.48\textwidth}
        \centering
        \includegraphics[width=\textwidth]{hpc2n_epsilon_cumulative_makespan.png}
        \caption{HPC2N -- $\epsilon$-Greedy}
        \label{fig:makespan_hpc2n_epsilon}
    \end{subfigure}
    \caption{Cumulative makespan trajectories across MAB variants on the HPC2N workload.}
    \label{fig:makespan_hpc2n_all}
\end{figure}

\begin{figure}[htbp]
    \centering
    \begin{subfigure}[b]{0.48\textwidth}
        \centering
        \includegraphics[width=\textwidth]{nasa_ucb_cumulative_makespan.png}
        \caption{NASA -- UCB}
        \label{fig:makespan_nasa_ucb}
    \end{subfigure}
    \hfill
    \begin{subfigure}[b]{0.48\textwidth}
        \centering
        \includegraphics[width=\textwidth]{nasa_discounted_ucb_cumulative_makespan.png}
        \caption{NASA -- Discounted UCB}
        \label{fig:makespan_nasa_discounted}
    \end{subfigure}
    \vskip\baselineskip
    \begin{subfigure}[b]{0.48\textwidth}
        \centering
        \includegraphics[width=\textwidth]{nasa_thompson_cumulative_makespan.png}
        \caption{NASA -- Thompson Sampling}
        \label{fig:makespan_nasa_thompson}
    \end{subfigure}
    \hfill
    \begin{subfigure}[b]{0.48\textwidth}
        \centering
        \includegraphics[width=\textwidth]{nasa_epsilon_cumulative_makespan.png}
        \caption{NASA -- $\epsilon$-Greedy}
        \label{fig:makespan_nasa_epsilon}
    \end{subfigure}
    \caption{Cumulative makespan trajectories across MAB variants on the NASA Ames workload.}
    \label{fig:makespan_nasa_all}
\end{figure}

The primary empirical insights include:
\begin{itemize}
    \item \textbf{Mitigation of Baseline Degradation:} Static \texttt{min\_min} exhibits severe performance degradation on NASA Ames (Figures~\ref{fig:makespan_nasa_ucb}--\ref{fig:makespan_nasa_epsilon}), accumulating over $1.6 \times 10^9$ seconds in makespan by round 100 due to severe task queueing. All MAB controllers avoid this failure mode by moving away from \texttt{min\_min} as early into the batch as possible.
    \item \textbf{Performance Superiority on HPC2N:} On HPC2N (Figure~\ref{fig:makespan_hpc2n_all}), standard UCB, Discounted UCB, Thompson Sampling, and $\epsilon$-Greedy all show a $59.5\%$--$71.4\%$ improvement over \texttt{min\_min} and constantly outperforms the static list schedulers.
    \item \textbf{Exploration Penalty under Thompson Sampling:} On NASA Ames (Figure~\ref{fig:makespan_nasa_thompson}), Thompson Sampling accumulates $8.0 \times 10^8$ seconds makespan, whereas $\epsilon$-Greedy (Figure~\ref{fig:makespan_nasa_epsilon}) terminates at $1.7 \times 10^8$ seconds, which shows that posterior sampling can lead to unnecessary exploration overhead in fairly predictable workloads.
\end{itemize}

% \subsubsection{Multi-Objective Radar Profile Analysis}
% Figures~\ref{fig:radar_hpc2n_all} and~\ref{fig:radar_nasa_all} compare normalized multi-objective trade-offs (Makespan, Cost, Carbon Footprint, and VM Utilization) on HPC2N and NASA Ames traces.

% \begin{figure}[htbp]
%     \centering
%     \begin{subfigure}[b]{0.48\textwidth}
%         \centering
%         \includegraphics[width=\textwidth]{hpc2n_ucb_radar_multimetric.png}
%         \caption{HPC2N -- UCB}
%         \label{fig:radar_hpc2n_ucb}
%     \end{subfigure}
%     \hfill
%     \begin{subfigure}[b]{0.48\textwidth}
%         \centering
%         \includegraphics[width=\textwidth]{hpc2n_discounted_ucb_radar_multimetric.png}
%         \caption{HPC2N -- Discounted UCB}
%         \label{fig:radar_hpc2n_discounted}
%     \end{subfigure}
%     \vskip\baselineskip
%     \begin{subfigure}[b]{0.48\textwidth}
%         \centering
%         \includegraphics[width=\textwidth]{hpc2n_thompson_radar_multimetric.png}
%         \caption{HPC2N -- Thompson Sampling}
%         \label{fig:radar_hpc2n_thompson}
%     \end{subfigure}
%     \hfill
%     \begin{subfigure}[b]{0.48\textwidth}
%         \centering
%         \includegraphics[width=\textwidth]{hpc2n_epsilon_radar_multimetric.png}
%         \caption{HPC2N -- $\epsilon$-Greedy}
%         \label{fig:radar_hpc2n_epsilon}
%     \end{subfigure}
%     \caption{Normalized multi-objective performance trade-offs on the HPC2N trace.}
%     \label{fig:radar_hpc2n_all}
% \end{figure}

% \begin{figure}[htbp]
%     \centering
%     \begin{subfigure}[b]{0.48\textwidth}
%         \centering
%         \includegraphics[width=\textwidth]{nasa_ucb_radar_multimetric.png}
%         \caption{NASA -- UCB}
%         \label{fig:radar_nasa_ucb}
%     \end{subfigure}
%     \hfill
%     \begin{subfigure}[b]{0.48\textwidth}
%         \centering
%         \includegraphics[width=\textwidth]{nasa_discounted_ucb_radar_multimetric.png}
%         \caption{NASA -- Discounted UCB}
%         \label{fig:radar_nasa_discounted}
%     \end{subfigure}
%     \vskip\baselineskip
%     \begin{subfigure}[b]{0.48\textwidth}
%         \centering
%         \includegraphics[width=\textwidth]{nasa_thompson_radar_multimetric.png}
%         \caption{NASA -- Thompson Sampling}
%         \label{fig:radar_nasa_thompson}
%     \end{subfigure}
%     \hfill
%     \begin{subfigure}[b]{0.48\textwidth}
%         \centering
%         \includegraphics[width=\textwidth]{nasa_epsilon_radar_multimetric.png}
%         \caption{NASA -- $\epsilon$-Greedy}
%         \label{fig:radar_nasa_epsilon}
%     \end{subfigure}
%     \caption{Normalized multi-objective performance trade-offs on the NASA Ames trace.}
%     \label{fig:radar_nasa_all}
% \end{figure}

% The radar profiles reveal structural trade-offs across algorithm classes:
% \begin{itemize}
%     \item \textbf{Operational Efficiency vs. Search Overhead:} Standalone \texttt{hghhc} achieves optimal internal schedule compaction but exhibits near-zero normalized scores ($< 0.1$) across Cost, Utilization, and Carbon axes due to its high metaheuristic execution overhead.
%     \item \textbf{Balanced Trade-off Envelope:} On NASA Ames (Figure~\ref{fig:radar_nasa_all}), the MAB hyper-heuristics achieve complete outer-bound dominance (normalized scores $> 0.95$ across all four axes), vastly outperforming single heuristics (\texttt{heft} and \texttt{peft} at $\approx 0.52$). On HPC2N (Figure~\ref{fig:radar_hpc2n_all}), MAB formulations maintain balanced scores between $0.65$ and $0.72$, ensuring sustainable resource utilization without incurring single-heuristic vulnerabilities.
% \end{itemize}

\subsubsection{Cross-Algorithm Selection Profile Heatmaps}
Figures~\ref{fig:heatmap_hpc2n} and~\ref{fig:heatmap_nasa} show how frequently each low-level heuristic was selected across all four MAB controllers on both the HPC2N and NASA workload traces.

\begin{figure}[htbp]
    \centering
    \begin{subfigure}[b]{0.48\textwidth}
        \centering
        \includegraphics[width=\textwidth]{heatmap_selection_hpc2n.png}
        \caption{HPC2N Trace Selection Profile}
        \label{fig:heatmap_hpc2n}
    \end{subfigure}
    \hfill
    \begin{subfigure}[b]{0.48\textwidth}
        \centering
        \includegraphics[width=\textwidth]{heatmap_selection_nasa.png}
        \caption{NASA Ames Trace Selection Profile}
        \label{fig:heatmap_nasa}
    \end{subfigure}
    \caption{Cross-algorithm low-level heuristic selection frequency profiles (\%) across MAB variants for the HPC2N and NASA Ames workload traces.}
    \label{fig:heatmaps_both}
\end{figure}

Key selection tendencies across workload domains include:
\begin{enumerate}
    \item \textbf{Primary Arm Preference (HPC2N):} On the HPC2N trace (Figure~\ref{fig:heatmap_hpc2n}), \texttt{hghhc} dominates selection across all MAB algoorithms due to the fact that it can efficiently schedule the tasks into available virtual machines.
    \item \textbf{Workload-Driven Policy Shift (NASA Ames):} On the NASA Ames trace (Figure~\ref{fig:heatmap_nasa}), the selection density shifts dramatically away from pure metaheuristics toward list-based and priority-driven heuristics. \texttt{peft} and \texttt{sufferage} emerge as preferred arms under UCB and Discounted UCB (accounting for $\approx 38.0\%$ and $26.0\%$ of total selections, respectively), while \texttt{round\_robin} and \texttt{max\_min} are sparsely engaged during high-concurrency rounds.
    \item \textbf{Filtering Inefficient Arms:} Inefficient or high-overhead heuristics are suppressed, receiving $\le 4.0\%$ of total choices across all MAB engines.
\end{enumerate}


% \subsection{Contextualization with Prior Research}
% Our empirical findings support and expand upon key themes in recent literature:
% \begin{itemize}
%     \item \textbf{Adaptive Heuristic Selection:} Consistent with Hosseini-Shirvani et al., fixed heuristics suffer from structural rigidity across non-stationary traces. Our MAB framework overcomes this by dynamically selecting high-performing heuristics based on real-time feedback.
%     \item \textbf{Non-Stationary Regret Control:} In alignment with non-stationary bandit theory \citep{Singh2024Leveraging, DeCarvalho2025Toward}, exponential history decay in Discounted UCB prevents Q-value stagnation, providing smoother makespan growth compared to unweighted Thompson Sampling.
%     \item \textbf{Resource Overhead Management:} While standalone metaheuristics like \texttt{hghhc} provide high static schedule quality, their unconstrained search incurs high energy costs. Operating at the hyper-heuristic level achieves strong schedule quality while managing overall operational overhead.
% \end{itemize}

\subsection{Critical Appraisal and Methodological Limitations}
To ensure complete transparency, we identify three methodological limitations of this study:

\subsubsection{Hyperparameter Tuning and Exploration Bounds}
The exploration factor $c = \sqrt{2}$ in UCB and discount factor $\gamma = 0.95$ in Discounted UCB were used as a global constant, which can lead to longer exploration phases than needed, which may favour sub-optimal heuristics. Finding a way to dynamically set these values will ensure that the scheduler maintains a decent balance between exploration and exploitation.

\subsubsection{Simulation Idealization vs. Real-World Cloud Dynamics}

For the purpose of this project, we make use of CloudSim, which executes tasks deterministically. But this will never be the case in real-world scenarios, as there will be network jitter, Vm cold starts and hypervisor resource contention, which will introduce some level of variance into heuristic execution times.

To bridge the gap between deterministic simulation and physical cloud infrastructure, three architectural extensions can be made: Asynchronous Policy Updates that transition the Upper-Level Multi-Armed Bandit (MAB) engine from discrete batch execution to an asynchronous event-driven loop that updates reward metrics continuously as telemetry streams arrive from live cloud nodes, adaptive reward smoothing that implements an exponential moving average on raw runtime feedback to prevent transient network anomalies from corrupting historical reward distributions and Fallback \& Safety Bounds by Setting hard SLA timeout thresholds that force a deterministic baseline heuristic (e.g., HEFT) if an exploitation arm experiences unexpected physical delay.

To bridge this gap, we can make three extensions:

\begin{itemize}
    \item \textbf{Asynchronous Policy Update:} This is going to help transition the decision plane from a discrete batch execution model to an asynchronous event-driven loop that updates reward metrics continuously.
    \item \textbf{Adaptive Reward Smoothing:} An exponential moving average on feedback to prevent anomalies from corrupting historical reward distributions.
    \item \textbf{Fallback and Safety Bounds:} Hard SLA thresholds should be set that forces a baseline heuristic, if an exploitation arm experiences any delay.
\end{itemize}
Deploying the current system directly onto live cloud infrastructure without these adaptation mechanisms yields predictable performance trade-offs such as an exploration phase delay, a robustness to non-stationary shift and suboptimal fine-grained decisions.


\subsubsection{Coarse Granularity of Decision Horizons}
Evaluating performance at batch boundaries hides task-level variations within a batch. Extremely long tasks can delay reward feedback, temporarily slowing policy adaptation.

\subsection{Suggested Modifications and Design Improvements}
To address these limitations, we propose three enhancements:
\begin{enumerate}
    \item \textbf{Adaptive Memory Decay ($\gamma_k$):} We can opt to adjust the discount factor for the older actions dynamically, based on the workload volatility. 
    \item \textbf{Contextual Bandits:} Incorporate state vectors to enable state-aware heuristic selection.
    \item \textbf{Testbed Deployment:} We can evaluate the scheduler on an actual Kubernetes cluster to assess the performance under real-world circumstance.
\end{enumerate}